\documentclass{report}
\usepackage{amsmath,amssymb}
\setlength{\parindent}{0mm}
\setlength{\parskip}{1em}
\begin{document}
\begin{center}
\rule{6in}{1pt} \
{\large Matthias Bollhoefer \\
{\bf Sparse Inverse Covariance Estimation Using Sparse Matrices}}

TU Braunschweig \\ Institute for Computational Mathematics \\ D-38106 Braunschweig \\ Germany
\\
{\tt m.bollhoefer@tu-bs.de}\\
Olaf Schenk\\
Fabio Verbosio\end{center}

A well-known task in statistics is the estimation of the (inverse)
covariance matrix when only certain samples of random variables are
given. In particular, big data problems often lead to a high-dimensional
regime where the number of samples is significant less than the
dimension. To compute an inverse covariance often requires additional
constraints such as sparsity. This is why one usually reguluarizes the
log-likelihood function with the 1-norm of the inverse covariance matrix.
In [1], a promising quadratic approximation called QUIC algorithm is
presented to estimate the sparse inverse covariance matrix for dense
matrices. This method has been generalized in [2] using the notion of
hierarchical matrices. In this talk we present a sparse version of this
algorithm which uses up-to-date sparse matrix computations to deal with
large-scale data and to handle sparse matrices and their approximate
inverses.

[1] Cho-Jui Hsieh, M�ty�s A. Sustik, Inderjit S. Dhillon, Pradeep
Ravikumar. Sparse Inverse Covariance Matrix Estimation Using Quadratic
Approximation, Advances in Neural Information Processing Systems, vol.
24, 2011.

[2] Jonas Ballani, Daniel Kressner.
Sparse Inverse Covariance Estimation with Hierarchical Matrices.
Technical report, EPFL, October 2014.


\end{document}
